\begin{acknowledgements}

岁月的车轮悄然转过，两载有余的硕士研究生生涯即将画上句号。回望这段旅程，从初试录取的欣喜、踏入科大校园的憧憬，到投身企业实习的历练，再到如今论文定稿的如释重负，个中收获不仅体现在学术能力的提升，更在于心智的成长。值此论文完稿之际，谨向所有给予我指导、支持与陪伴的人致以最诚挚的谢意。

首先，衷心感谢我的导师赵功名教授。本论文从选题构思、开题论证到最终定稿，每一步都凝结着赵老师的心血。赵老师严谨治学的态度与敏锐的学术洞察力令我获益匪浅。特别难忘在开题阶段，赵老师通过电话针对论文题目的准确性给予了关键指导，使研究重点得以聚焦；在论文撰写过程中，赵老师更是多次审阅稿件，从逻辑架构到遣词造句均提出了详尽的修改意见，确保了本研究的严密性与规范性。在此，向赵老师表达崇高的敬意。同时，也要感谢张钟高级工程师在项目实践中给予的宝贵建议。

其次，感谢中国科学技术大学软件学院及全体授课教师。学院设置的专业课程体系兼具前沿性与实用性，为我的科研工作打下了坚实的理论基础。同时，感谢学院搭建的实习双选平台，使我有机会进入业界领军企业进行工程实践，在真实场景中提升了解决复杂问题的能力。

感谢我的父母。考研备考至今，父母始终是我最坚强的后盾。正是因为他们在经济上的无私支持与精神上的默默守护，使我能够免于生活琐事的后顾之忧，全身心地投入到学业与研究之中。这份恩情，厚重如山，没齿难忘。

感谢与我并肩作战的同学们。无论是在课程研讨中的思维碰撞，还是在实习期间的相互勉励，大家的陪伴使这段旅程充满了温暖。在论文撰写期间，我们针对学术难题进行的探讨与互助审校，极大地拓宽了我的视野。这段纯粹的同窗情谊是我硕士生涯中最珍贵的财富。

最后，感谢母校中国科学技术大学。“红专并进，理实交融”的校训精神已深深融入我的血液。能在这所享有盛誉的学府深造，是我一生的荣幸与骄傲。离别之际，唯愿母校再创辉煌，愿恩师教祺，愿同窗前程似锦。

\end{acknowledgements}
